\chapter{MARTINGALE THEORY}

\section{Discrete Time Martingales}
\textit{September 16th.}

\begin{definition}
    A sequence of random variables $\{Z_{n}\}_{n \geq 1}$ is a \eax{martingale} if
    \begin{enumerate}
        \item $\E[\abs{Z_{n}}] < \infty$ for all $n \geq 1$,
        \item $\E[Z_{n} \mid Z_{n-1} = z_{n-1}, \ldots, Z_{1} = z_{1}] = z_{n-1}$.
    \end{enumerate}
\end{definition}
To understand the second condition, let us focus on the following: with $(\Omega,\cF,\Prob)$ as a probability space, and $Y,X:\Omega \to S$ be discrete random variables with $S$ a countable set, then with $y \in Y(\Omega)$, we have
\begin{align}
    \E[X \mid Y = y] = \sum_{k \in S} k \Prob(X = k \mid Y = y)
\end{align}
For multiple discrete random variables $Y_{1},Y_{2},\ldots,Y_{n}$, we have
\begin{align}
    \E[X \mid Y_{1}=y_{1},\ldots,Y_{n}=y_{n}] = \sum_{k \in S} k \Prob(X=k \mid Y_{1}=y_{1},\ldots,Y_{n}=y_{n})
\end{align}
for $y_{i} \in Y_{i}(\Omega)$.

\begin{example}
    Let $\{S_{n}\}_{n \geq 1}$ be the simple symmetric random walk on $\Z$, with $S_{0} = 0$, and let $s_{1},\ldots,s_{n-1} \in Z$ be such that $\Prob(S_{n-1}=s_{n-1},\ldots,S_{1}=s_{1}) > 0$. As a formal or imprecise calculation, we have
    \begin{align}
        \E[S_{n} \mid S_{n-1}=s_{n-1},\ldots,S_{1}=s_{1}] &= \sum_{k \in \Z} k \Prob(S_{n}=k \mid S_{n-1}=s_{n-1},\ldots,S_{1}=s_{1}) \notag \\
        &= \sum_{k =-n}^{n} k \Prob(S_{n-1}+X_{n}=k \mid S_{n-1}=s_{n-1},\ldots,S_{1}=s_{1}) \notag \\
        &= \sum_{k =-n}^{n} k \Prob(X_{n}=k-s_{n-1} \mid S_{n-1}=s_{n-1},\ldots,S_{1}=s_{1}) \notag \\
        &= \sum_{k =-n}^{n} k \Prob(X_{n}=k-s_{n-1}) \notag \\
        &= (s_{n-1}-1)\Prob(X_{n}=1) + (s_{n-1}+1)\Prob(X_{n}=-1) = s_{n-1}.
    \end{align}
    Also, $\E[\abs{S_{n}}] < \infty$ for all $n \geq 1$. If instead, we consider the simple asymmetric random walk on $\Z$ (that is, $\Prob(X_{n}=1) = p = 1-\Prob(X_{n}=-1)$ and $p \neq \frac{1}{2}$), then we have
    \begin{align}
        \E[S_{n} \mid S_{n-1}=s_{n-1},\ldots,S_{1}=s_{1}] &= (s_{n-1}-1)p + (s_{n-1}+1)(1-p) = s_{n-1} + 1 - 2p \neq s_{n-1}.
    \end{align}
    Hence, the simple symmetric random walk is a martingale, while the simple asymmetric random walk is not.
\end{example}

\begin{example}
    Let $X_{n}$ be a sequence of independent and identically distributed random variables with $\Prob(X_{n} = 2) = \frac{1}{2} = \Prob(X_{n} = 0)$, and then define $Z_{n} = \prod_{i=1}^{n} X_{i}$. Then, we have
    \begin{align}
        \E[\abs{Z_{n}}] = \prod_{k=1}^{n} \E[\abs{X_{k}}] = \prod_{k=1}^{n} 1 = 1 < \infty
    \end{align}
    for all $n \geq 1$. Now, for $z_{1},\ldots,z_{n-1} \in \{0,2,4,\ldots,2^{n-1}\}$ such that $\Prob(Z_{n-1}=z_{n-1},\ldots,Z_{1}=z_{1}) > 0$, we have
    \begin{align}
        \E[Z_{n} \mid Z_{n-1}=z_{n-1},\ldots,Z_{1}=z_{1}] &= \E[X_{n} Z_{n-1} \mid Z_{n-1}=z_{n-1},\ldots,Z_{1}=z_{1}] \notag \\
        &= z_{n-1} \E[X_{n} \mid Z_{n-1}=z_{n-1},\ldots,Z_{1}=z_{1}] \notag \\
        &= z_{n-1} \E[X_{n}] = z_{n-1}.
    \end{align}
    Hence, $\{Z_{n}\}_{n \geq 1}$ is a martingale.
\end{example}

Suppose $\{Z_{n}\}_{n \geq 1}$ is a Markov chain. A consequence of the Markov property,
\begin{align}
    \Prob(Z_{n}=z_{n},\ldots,Z_{0}=z_{0}) = \mu(\{z_{0}\}) \prod_{k=1}^{n} p_{z_{k-1},z_{k}},
\end{align}
is that the past and the future are independent given the present. A martingale is less rigid than a Markov chain, saying that the expectation of the future given the present and the past is equal to the present. Moreover, Martingales are easier to construct and verify than Markov chains, and several limit results follow too.

\begin{example}
    Let $Z_{1}$ be the random variable defined by $\Prob(Z_{1}=1) = \frac{1}{2} = \Prob(Z_{1}=0)$. If $\{X_{n}\}_{n \geq 1}$ is a sequence of independent and identically distributed random variables with $\Prob(X_{n}=1) = \frac{1}{2} = \Prob(X_{n}=-1)$, then define $Z_{n} = Z_{n-1} + X_{n}Z_{1}$ for $n \geq 2$. Then, verify that $\{Z_{n}\}_{n \geq 1}$ is a Markov chain, but not a martingale.
\end{example}

\noindent From here on, we will interpet $\E[Z_{n} \mid z_{n-1},\ldots,z_{1}]$ to be $\E[Z_{n} \mid Z_{n-1}=z_{n-1},\ldots,Z_{1}=z_{1}]$.

Let $\{Z_{n}\}_{n \geq 1}$ be a collection of discrete random variables on $(\Omega, \cF, \Prob)$. Then we know that
\begin{align}
    \E[Z_{n} \mid Z_{n-1}=z_{n-1},\ldots,Z_{1}=z_{1}] = \sum_{k \in Z_{n}(\Omega)} k \Prob(Z_{n}=k \mid Z_{n-1}=z_{n-1},\ldots,Z_{1}=z_{1}).
\end{align}
Note that both sides of the above equation is a function of $z_{1},\ldots,z_{n-1}$, say, $f(z_{1},\ldots,z_{n-1})$. If $\{Z_{n}\}_{n \geq 1}$ were a martingale, then the equation above implies that $f(z_{1},\ldots,z_{n-1}) = z_{n-1}$, which is a function of $z_{n-1}$ alone. Thus, we define $Y_{n-1} \defeq f(z_{1},\ldots,z_{n-1})$, which we interpret as $\E[Z_{n} \mid z_{n-1},\ldots,z_{1}]$.

\begin{remark}
    We recall two important facts.
    \begin{enumerate}
        \item If $Y$ and $Z$ are discrete random variables, and $X = h(Y,Z)$ for some function $h:\R^{2}\to \R$, then $\E[X \mid Y=y,Z=z] = h(y,z)$.
        \item If $X$ and $Y$ are independent discrete random variables, then $\E[X \mid Y=y] = \E[X] \equiv \E[X \mid Y]$.
    \end{enumerate}
\end{remark}

\begin{example}
    Let $\{X_{n}\}_{n \geq 1}$ be a sequence of independent and identically distributed random variables with $\Prob(X_{n}=1) = \frac{1}{2} = \Prob(X_{n}=-1)$. Define $S_{n} = \sum_{i=1}^{n} X_{i}$ for $n \geq 1$, and $S_{0} = 0$. Now, define $\tilde{Z}_{n} = e^{rS_{n}}$ for all $n \geq 1$, where $r \in \R$ is a constant. Then, we have
    \begin{align}
        \E[\abs{\tilde{Z}_{n}}] = \E[e^{rS_{n}}] \leq c < \infty
    \end{align}
    as $\abs{S_{n}} \leq n$. Now, notice that $e^{rS_{n}} = e^{r\sum_{i=1}^{n} X_{i}} = e^{r\sum_{i=1}^{n} X_{i}}$. But with $\tilde{Y}_{i} = e^{rX_{i}}$, we have $\tilde{Y}_{1},\ldots,\tilde{Y}_{n}$ are independent and identically distributed random variables, and $\tilde{Z}_{n} = \prod_{i=1}^{n} \tilde{Y}_{i}$. This gives $\E[X_{i}] = \frac{e^{r}+e^{-r}}{2} \neq 1$ unless $r = 0$. Hence, $\tilde{Z}_{n}$ is not a martingale. To extract a martingale out of this, define
    \begin{align}
        q(r) = \log \E[\tilde{Y}_{i}] = \log \left( \frac{e^{r}+e^{-r}}{2} \right),
    \end{align}
    and let $Y_{i} = e^{rX_{i}-q(r)}$. Then, we have that $Y_{i}$'s are independent and identically distributed random variables with $\E[Y_{i}] = 1$. Now, define $Z_{n} = \prod_{i=1}^{n} Y_{i} = e^{rS_{n}-nq(r)}$. Then, we have, via the same argument as in \textbf{Example 3.3}, that $\{Z_{n}\}_{n \geq 1}$ is a (product) martingale. More generally, if $\{X_{n}\}_{n \geq 1}$ is a sequence of independent and identically distributed random variables, and $S_{n} = \sum_{i=1}^{n} X_{i}$, then $\{Z_{n}\}_{n \geq 1}$ defined by $Z_{n} = e^{rS_{n}-nq(r)}$ is a martingale, where $q(r) = \log \E[e^{rX_{1}}]$.
\end{example}

\begin{lemma}
    Let $\{Z_{n}\}_{n \geq 1}$ be a martingale. Then, for any $n \geq 1$, we have
    \begin{align}
        \E[Z_{n} \mid z_{i},z_{i-1},\ldots,z_{1}] = z_{i}
    \end{align}
    for all $1 \leq i \leq n$.
\end{lemma}
\begin{proof}
    We make use of the tower property of conditional expectation, which states that for any random variables $X,Y,Z$, we have
    \begin{align}
        \E[\E[X \mid Y,Z] \mid Y] = \E[X \mid Y].
    \end{align}
    We prove via induction. The base case is $n=1$, which is trivial. Now, assume the statement holds for $n = m$, and let $n = m+1$. Then,
    \begin{align}
        \E[Z_{m+1} \mid z_{i},z_{i-1},\ldots,z_{1}] &= \E[\E[Z_{m+1} \mid z_{m},z_{m-1},\ldots,z_{1}] \mid z_{i},z_{i-1},\ldots,z_{1}] \notag \\
        &= \E[Z_{m} \mid z_{i},z_{i-1},\ldots,z_{1}] = z_{i},
    \end{align}
    where the second equality follows from the martingale property, and the third equality follows from the induction hypothesis. Hence, by induction, the statement holds for all $n \geq 1$.
\end{proof}
If we apply the above lemma to $i = 1$, then
\begin{align}
    \E[Z_{n} \mid Z_{1}] = Z_{1} \implies \E[Z_{n}] = \E[\E[Z_{n} \mid Z_{1}]] = \E[Z_{1}],
\end{align}
for all $n \geq 1$. Hence, the expectation of a martingale is constant over time.

\begin{remark}
    We show the tower property of conditional expectation for discrete random variables. Let $X,Y,Z$ be discrete random variables, and let $y \in Y(\Omega)$ and $z \in Z(\Omega)$. Let
    \begin{align}
        h(y,z) = \E[X \mid Y=y,Z=z] = \sum_{k \in X(\Omega)} k \Prob(X=k \mid Y=y,Z=z).
    \end{align}
    Then,
    \begin{align}
        \E[h(Y,Z) \mid Y=y_{0}] &= \sum_{y \in Y(\Omega)} \sum_{z \in Z(\Omega)} h(y,z) \Prob(Y=y,Z=z \mid Y=y_{0}) = \sum_{z \in Z(\Omega)} h(y_{0},z) \Prob(Z=z \mid Y=y_{0}) \notag \\
        &= \sum_{z \in Z(\Omega)} \left( \sum_{k \in X(\Omega)} k \Prob(X=k \mid Y=y_{0},Z=z) \right) \Prob(Z=z \mid Y = y_{0}) \notag \\
        &= \sum_{z \in Z(\Omega)} \sum_{k \in X(\Omega)} k \cdot \frac{\Prob(X=k,Y=y_{0},Z=z)}{\Prob(Y=y_{0})} \notag \\
        &= \frac{1}{\Prob(Y=y_{0})}\sum_{z \in Z(\Omega)} \sum_{k \in X(\Omega)} k \Prob(X=k,Y=y_{0},Z=z) \notag \\
        &= \frac{1}{\Prob(Y=y_{0})}\sum_{k \in X(\Omega)} k \Prob(X=k,Y=y_{0}) = \E[X \mid Y=y_{0}].
    \end{align}
    Hence, $\E[h(Y,Z) \mid Y=y_{0}] = \E[X \mid Y=y_{0}]$, which is the tower property of conditional expectation.
\end{remark}